\documentclass[12pt,a4paper]{article}

\usepackage{amsmath,amssymb,amsthm}
\usepackage{geometry}
\usepackage{booktabs}
\usepackage{array}
\usepackage{hyperref}
\usepackage{setspace}
\usepackage{parskip}

\geometry{margin=1in}
\onehalfspacing

\newtheorem{proposition}{Proposition}
\newtheorem{corollary}{Corollary}
\newtheorem{definition}{Definition}
\newtheorem{remark}{Remark}

\title{\textbf{Wage Inequality, Upskilling, and Long-Run Growth:\\
A Three-Player Model}\\[0.8em]
\large Working Paper}
\author{}
\date{\today}

\begin{document}

\maketitle

\begin{abstract}
We develop a tractable three-player model connecting wage inequality,
worker upskilling decisions, and firm asset accumulation. A single
parameter $\mu \in (0,\tfrac{1}{2})$ --- the income share of the
lower-ranked worker group --- governs the entire distribution of wages.
Workers allocate income between upskilling investment and participation
in a zero-sum displacement lottery. Preferences are modelled using
linear utility augmented by a signed curvature parameter $\theta_s$
per worker: $\theta_L > 0$ for the lower-income worker (risk-seeking)
and $\theta_H < 0$ for the higher-income worker (risk-averse). This
two-parameter specification replaces the full Friedman-Savage machinery
while preserving the essential risk asymmetry. All equilibrium
quantities --- upskilling fractions, lottery stakes, firm asset
accumulation, and steady-state wage split --- remain in closed form.
We distinguish two upskilling measures: an individual measure
$K^*_{\mathrm{ind}} = 2\kappa^{\min}(\mu)$ independent of population
structure, and an aggregate measure
$K^*_{\mathrm{agg}} = (W+1)\kappa^{\min}(\mu)$ weighted by the
population mass ratio $W = i^*/(N-i^*)$, which enters firm growth
through the upskilling index $\Phi(\mu)$. We address two questions:
how do workers respond to wage inequality $\mu$, and what constraints
does $\mu$ place on long-term growth?
\end{abstract}

\tableofcontents
\newpage

%----------------------------------------------------------------------
\section{Introduction}
%----------------------------------------------------------------------

How should a firm set wages across skill groups to maximise long-run
asset growth? And how do workers in different income groups respond to
the wage structure they face? This paper provides a unified framework
with closed-form answers to both questions.

There are three players: a lower-income worker $L$, a higher-income
worker $H$, and a firm. The firm holds productive assets and pays wages
as a share of those assets. It chooses a wage-split parameter
$\mu \in (0, \tfrac{1}{2})$ each period --- the income share of the
lower-income group --- trading off the upskilling return from greater
equality against a convex restructuring cost. Workers receive wages
determined by $\mu$ and split income between upskilling investment and
participation in a zero-sum displacement lottery. A participation
constraint ensures each worker invests enough to maintain a subsistence
base income.

Preferences are modelled parsimoniously. The baseline is linear utility
$u(m) = m$, augmented by a single curvature correction per worker:
\begin{equation}
  u_s(m) = m + \frac{\theta_s}{2}(m - \bar{b})^2
  \label{eq:utility_intro}
\end{equation}
where $\bar{b}$ is the subsistence income and $\theta_s = u_s''(\bar{b})$
is the signed curvature at subsistence. Setting $\theta_L > 0$ makes
worker $L$ locally risk-seeking (convex utility) and $\theta_H < 0$
makes worker $H$ locally risk-averse (concave utility). This replaces
the three-parameter Friedman-Savage specification $(\alpha_u, \beta_u,
m^*)$ with two parameters $(\theta_L, \theta_H)$, while retaining the
key qualitative asymmetry. Crucially, these parameters affect only the
welfare evaluation of lottery participation, not the equilibrium
allocation quantities, which remain unchanged from the linear case.

A central structural result concerns the role of the population mass
ratio $W = i^*/(N-i^*)$, where $i^*$ is the class barrier separating
lower- and upper-income workers in the underlying $N$-worker model.
At the individual level, each representative worker's upskilling
investment $\kappa^{\min}(\mu)$ is determined solely by the
participation constraint, with no dependence on $W$. At the
economy-wide level, however, the population-weighted aggregate
$(W+1)\kappa^{\min}(\mu)$ does depend on $W$ and is the quantity the
firm internalises through its upskilling index $\Phi(\mu)$. This
distinction --- $W$ enters growth through aggregation, not through
individual incentives --- is a main structural result of the model.

%----------------------------------------------------------------------
\section{Model Setup}
%----------------------------------------------------------------------

\subsection{Players and Incomes}

There are three strategic players: worker $L$, worker $H$, and a firm.
The firm holds asset stock $A_t$ and distributes total income $R$ to
workers. It chooses the income-split parameter $\mu \in (0, \tfrac{1}{2})$:
\begin{equation}
  M_L = R\mu, \qquad M_H = R(1-\mu)
  \label{eq:incomes}
\end{equation}
Since $\mu < \tfrac{1}{2}$, worker $H$ always earns more than worker
$L$. Higher $\mu$ means greater wage equality; lower $\mu$ means
greater inequality. The income ratio $M_H/M_L = (1-\mu)/\mu$ is
strictly decreasing in $\mu$.

\begin{definition}[Wage inequality parameter]
$\mu \in (0, \tfrac{1}{2})$ is the income share of worker $L$,
chosen by the firm each period. It is the single measure of wage
inequality in the model.
\end{definition}

In the underlying $N$-worker model with a Pareto income ladder
$m_i = x_0(1 - i/(N+1))^{-1/\alpha}$ and a class barrier at rank
$i^*$, the two-player reduction yields $\mu$ as the lower hump's
share of total representative income:
\begin{equation}
  \mu \approx \frac{\displaystyle\int_0^{i^*/N}
    x_0(1-s)^{-1/\alpha}\,ds}
  {\displaystyle\int_0^{1}
    x_0(1-s)^{-1/\alpha}\,ds}
  \label{eq:mu_micro}
\end{equation}
In the macroeconomic model $\mu$ is taken as the firm's policy
instrument rather than a fixed structural outcome.

\subsection{The Population Mass Ratio $W$}

The population mass ratio is:
\begin{equation}
  W = \frac{i^*}{N - i^*} \geq 1
  \label{eq:W_def}
\end{equation}
Under the Pareto distribution the lower group is generically more
numerous ($W \geq 1$). The parameter $W$ enters the model in two
distinct and separate ways:
\begin{enumerate}
  \item \textbf{Lottery pool}: $W$ amplifies worker $L$'s aggregate
        stake in the prize pot, reflecting the greater frequency of
        lower-rank displacement events per period.
  \item \textbf{Aggregate upskilling}: $W$ weights the lower group's
        upskilling in the population-weighted total $(W+1)\kappa^{\min}$,
        which the firm uses to compute the return on its wage bill.
\end{enumerate}
$W$ does \emph{not} affect any individual worker's budget or
participation constraint. Each representative worker invests
$\kappa^{\min}(\mu)$ regardless of $W$.

\subsection{Worker Budget and Allocation}

Each worker $s \in \{L, H\}$ allocates fraction $\lambda_s \in [0,1]$
of income $M_s$ to upskilling and the remainder to the lottery:
\begin{equation}
  \kappa_s = \lambda_s M_s \quad \text{(upskilling)},
  \qquad
  \ell_s = (1-\lambda_s)M_s \quad \text{(lottery stake)}
  \label{eq:budget}
\end{equation}
Upskilling produces output $f(\kappa_s) = \kappa_s^\gamma$,
$\gamma \in (0,1)$, with diminishing returns. The certain base income
after upskilling and human capital depreciation $\delta_w > 0$:
\begin{equation}
  b_s(\lambda_s) = \kappa_s + \kappa_s^\gamma - \delta_w
  \label{eq:base_income}
\end{equation}

\subsection{Two Upskilling Measures}

\begin{definition}[Individual upskilling total]
$K^*_{\mathrm{ind}} = \kappa_L^* + \kappa_H^* = 2\kappa^{\min}(\mu)$
is the sum of the two representative agents' equilibrium investments.
It is independent of $W$.
\end{definition}

\begin{definition}[Aggregate upskilling total]
$K^*_{\mathrm{agg}} = W\kappa_L^* + \kappa_H^* = (W+1)\kappa^{\min}(\mu)$
is the population-weighted sum. It enters the firm's asset accumulation
through the upskilling index $\Phi(\mu) = K^*_{\mathrm{agg}}/R$.
\end{definition}

The ratio $K^*_{\mathrm{agg}}/K^*_{\mathrm{ind}} = (W+1)/2 \geq 1$.
For $W > 1$, aggregate upskilling exceeds the two-agent sum because
the more numerous lower group amplifies the economy-wide human capital
base.

\subsection{Zero-Sum Displacement Lottery}

The lottery pool is:
\begin{equation}
  \Lambda = W(1-\lambda_L)M_L + (1-\lambda_H)M_H = W\ell_L + \ell_H
  \label{eq:pool}
\end{equation}
Each worker wins the pool with probability $\tfrac{1}{2}$. Outcome
incomes are:
\begin{align}
  M_L^{\mathrm{win}}  &= b_L + (1-\lambda_H)M_H \label{eq:MLwin} \\
  M_L^{\mathrm{lose}} &= b_L - W(1-\lambda_L)M_L \label{eq:MLlose} \\
  M_H^{\mathrm{win}}  &= b_H + W(1-\lambda_L)M_L \label{eq:MHwin} \\
  M_H^{\mathrm{lose}} &= b_H - (1-\lambda_H)M_H \label{eq:MHlose}
\end{align}
The lottery is strictly zero-sum: $\Lambda$ is redistributed, not
created. Only upskilling changes total income.

\subsection{Participation Constraint}

Each worker must maintain base income above a subsistence floor
$\bar{b}(\mu) = \bar{b}_0 R\mu$, proportional to the lower worker's
wage:
\begin{equation}
  b_s(\lambda_s) \geq \bar{b}(\mu)
  \iff
  \kappa_s + \kappa_s^\gamma \geq \bar{b}_0 R\mu + \delta_w
  \label{eq:pc}
\end{equation}
The minimum upskilling satisfying this constraint, $\kappa^{\min}(\mu)$,
solves:
\begin{equation}
  \kappa^{\min} + \left(\kappa^{\min}\right)^\gamma
  = \bar{b}_0 R\mu + \delta_w
  \label{eq:kappa_min}
\end{equation}
For $\gamma = 2$, adopted throughout for closed form:
\begin{equation}
  \kappa^{\min}(\mu)
  = \frac{-1 + \sqrt{1 + 4(\bar{b}_0 R\mu + \delta_w)}}{2}
  \label{eq:kappa_min_gamma2}
\end{equation}
This is strictly increasing and concave in $\mu$, with sensitivity:
\begin{equation}
  \frac{d\kappa^{\min}}{d\mu} \equiv \Psi(\mu)
  = \frac{\bar{b}_0 R}{1 + 2\kappa^{\min}(\mu)} > 0
  \label{eq:dkdmu}
\end{equation}

%----------------------------------------------------------------------
\section{Preferences: Linear Utility with Curvature Correction}
%----------------------------------------------------------------------

\subsection{The Two-Parameter Specification}

We use linear utility augmented by a quadratic curvature correction
evaluated at the subsistence income $\bar{b}$:
\begin{equation}
  u_s(m) = m + \frac{\theta_s}{2}(m - \bar{b})^2
  \label{eq:utility}
\end{equation}
where $\theta_s \in \mathbb{R}$ is the curvature parameter for worker
$s$. The derivatives are:
\begin{align}
  u_s'(m) &= 1 + \theta_s(m - \bar{b}) \label{eq:uprime} \\
  u_s''(m) &= \theta_s \label{eq:upprime}
\end{align}
The curvature $u_s'' = \theta_s$ is constant --- the simplest possible
departure from linearity. We impose:
\begin{equation}
  \theta_L > 0 \quad \text{(worker $L$: convex utility, risk-seeking)},
  \qquad
  \theta_H < 0 \quad \text{(worker $H$: concave utility, risk-averse)}
  \label{eq:theta_signs}
\end{equation}
This two-parameter specification directly encodes the Friedman-Savage
risk asymmetry without requiring an inflection point or a polynomial
of degree greater than two.

\begin{remark}
Equation~(\ref{eq:utility}) is equivalent to a second-order Taylor
expansion of any smooth utility function around $\bar{b}$, retaining
only the zeroth, first, and second-order terms. The parameter
$\theta_s = u_s''(\bar{b})$ is the Arrow-Pratt numerator evaluated at
subsistence. The Arrow-Pratt denominator $u_s'(\bar{b}) = 1$ under
this specification, giving absolute risk aversion
$A_s(\bar{b}) = -\theta_s$.
\end{remark}

\subsection{Risk Premia}

For a binary lottery paying $\bar{b} + x$ or $\bar{b} - y$ with equal
probability, the expected utility under specification~(\ref{eq:utility})
is:
\begin{equation}
  \frac{1}{2}u_s(\bar{b}+x) + \frac{1}{2}u_s(\bar{b}-y)
  = \bar{b} + \frac{x-y}{2}
    + \frac{\theta_s}{2}\cdot\frac{x^2+y^2}{2}
  \label{eq:EU_lottery}
\end{equation}
The first two terms are the linear expected value. The third term
$\frac{\theta_s}{4}(x^2 + y^2)$ is the risk premium correction.
Using the compact notation:
\begin{align}
  a(\mu) &= R\mu - \kappa^{\min}(\mu)
  && \text{(worker $L$ stake)} \label{eq:a} \\
  c(\mu) &= R(1-\mu) - \kappa^{\min}(\mu)
  && \text{(worker $H$ stake)} \label{eq:c} \\
  \Pi(\mu) &= Wa(\mu) + c(\mu)
  && \text{(total pool spread)} \label{eq:Pi} \\
  \epsilon_L(\mu) &= \tfrac{c(\mu) - Wa(\mu)}{2}
  && \text{(net expected transfer to $L$)} \label{eq:epsilonL} \\
  \epsilon_H(\mu) &= -\epsilon_L(\mu)
  && \text{(net expected transfer to $H$)} \label{eq:epsilonH}
\end{align}
the win and lose amounts for worker $s$ are
$x_s = \Pi_{-s}$ (prize received) and $y_s = W_s\ell_s$ (stake lost).
Substituting into equation~(\ref{eq:EU_lottery}):
\begin{equation}
  V_s = \bar{b} + \epsilon_s
        + \underbrace{\frac{\theta_s}{4}
        \!\left(\Pi_{-s}^2 + (W_s\ell_s)^2\right)}_{\tilde\rho_s}
  \label{eq:Vs_full}
\end{equation}
At the constrained equilibrium where $b_s = \bar{b}$,
$\Pi_L = c$, $\Pi_H = Wa$, $W_L\ell_L = Wa$, $W_H\ell_H = c$:
\begin{equation}
  \tilde\rho_L = \frac{\theta_L}{4}(c^2 + W^2a^2),
  \qquad
  \tilde\rho_H = \frac{\theta_H}{4}(W^2a^2 + c^2)
  \label{eq:rho_explicit}
\end{equation}
Writing $\Sigma^2 \equiv c^2 + W^2a^2$, the risk premia take the
symmetric form:
\begin{equation}
  \boxed{\tilde\rho_s = \frac{\theta_s}{4}\Sigma^2}
  \label{eq:rho_simple}
\end{equation}
Since $\theta_L > 0$: $\tilde\rho_L > 0$ --- worker $L$ gains welfare
from the lottery beyond expected income. Since $\theta_H < 0$:
$\tilde\rho_H < 0$ --- worker $H$ loses welfare from the lottery.

The net social risk premium is:
\begin{equation}
  \tilde\rho_L + \tilde\rho_H
  = \frac{\theta_L + \theta_H}{4}\,\Sigma^2
  \label{eq:social_premium}
\end{equation}
The lottery displacement mechanism is socially desirable if and only if:
\begin{equation}
  \theta_L + \theta_H > 0
  \label{eq:social_condition}
\end{equation}
The lower worker's risk-seeking gain outweighs the upper worker's
risk-aversion loss. This is a single inequality in two observable
parameters, requiring no reference to an inflection point.

\subsection{Equilibrium Utility}

At the constrained equilibrium ($b_s = \bar{b}$), the full utility
is:
\begin{equation}
  V_s = \underbrace{\bar{b}}_{\text{subsistence}}
      + \underbrace{\epsilon_s}_{\text{expected transfer}}
      + \underbrace{\frac{\theta_s}{4}\Sigma^2}_{\text{risk premium}}
  \label{eq:V_final}
\end{equation}
The three terms are cleanly separated. The curvature parameters
$(\theta_L, \theta_H)$ affect only the third term. They do not enter
the participation constraint, the optimal $\lambda_s^*$, the firm's
FOC, or any accumulation object. This clean separation is the main
advantage of the quadratic specification: risk preferences are welfare
relevant but allocation neutral.

%----------------------------------------------------------------------
\section{Derivation of Optimal $\lambda_s^*$}
%----------------------------------------------------------------------

\subsection{Why the FOC Has No Interior Solution}

Worker $s$ maximises $V_s(\lambda_s;\lambda_{-s})$ subject to
$b_s(\lambda_s) \geq \bar{b}$. Differentiating with respect to
$\lambda_s$:
\begin{equation}
  \frac{\partial V_s}{\partial\lambda_s}
  = \frac{\partial b_s}{\partial\lambda_s}
    + \frac{\partial\epsilon_s}{\partial\lambda_s}
    + \frac{\theta_s}{4}\frac{\partial\Sigma^2}{\partial\lambda_s}
  \label{eq:dV_full}
\end{equation}
Using $\partial b_s/\partial\lambda_s = M_s(1+\gamma\kappa_s^{\gamma-1})>0$,
$\partial\epsilon_s/\partial\lambda_s$ from the stake expressions, and
$\partial\Sigma^2/\partial\lambda_s$, each component is bounded and
finite. However, the dominant term is $\partial b_s/\partial\lambda_s
> 0$: raising $\lambda_s$ raises base income in both outcome states,
and for small $\theta_s$ this dominates. More precisely, taking the
full derivative and examining signs shows $\partial V_s/\partial\lambda_s
> 0$ over the relevant range, so $V_s$ is increasing in $\lambda_s$.
There is no interior zero of the FOC.

The economic reason is that raising $\lambda_s$ improves both the win
outcome ($b_s + \Pi_{-s}$ rises as $b_s$ rises) and the lose outcome
($b_s - W_s\ell_s$ rises as $b_s$ rises and $\ell_s$ falls). Both
states improve simultaneously, so the unconstrained optimum is the
corner $\lambda_s = 1$.

\subsection{The Binding Participation Constraint}

Since $V_s$ is increasing in $\lambda_s$ and the worker cannot set
$\lambda_s = 1$ without violating the subsistence floor
(at $\lambda_s = 1$, $\kappa_s = M_s$ but $b_s$ may exceed $\bar{b}$
--- the constraint binds from below as $\lambda_s$ is pushed up until
$b_s$ just reaches $\bar{b}$), the optimal solution is:
\begin{equation}
  b_s(\lambda_s^*) = \bar{b}
  \iff
  \kappa_s^* + (\kappa_s^*)^\gamma = \bar{b} + \delta_w
  \label{eq:binding_pc}
\end{equation}
This is precisely the equation defining $\kappa^{\min}$. Since
$\kappa_s^* = \lambda_s^* M_s = \lambda_s^* RM_s^0$:
\begin{equation}
  \lambda_s^* = \frac{\kappa^{\min}(\mu)}{RM_s^0}
  \label{eq:lambda_star_general}
\end{equation}
For worker $L$ ($M_L^0 = \mu$) and worker $H$ ($M_H^0 = 1-\mu$):
\begin{equation}
  \boxed{
    \lambda_L^*(\mu) = \frac{\kappa^{\min}(\mu)}{R\mu},
    \qquad
    \lambda_H^*(\mu) = \frac{\kappa^{\min}(\mu)}{R(1-\mu)}
  }
  \label{eq:lambda_star}
\end{equation}

\subsection{Ordering and Separation}

Since $R\mu < R(1-\mu)$, the ordering $\lambda_L^* > \lambda_H^*$
holds always. The lower-income worker invests a larger fraction in
upskilling because the same absolute floor $\kappa^{\min}$ represents
a larger share of a smaller budget. This is an income-scaling effect,
not a preference effect. The curvature parameters $(\theta_L, \theta_H)$
do not appear in equation~(\ref{eq:lambda_star}).

%----------------------------------------------------------------------
\section{Question 1: Worker Responses to Wage Inequality $\mu$}
\label{sec:workers}
%----------------------------------------------------------------------

\subsection{Upskilling Responses}

Differentiating equation~(\ref{eq:lambda_star}):
\begin{align}
  \frac{d\lambda_L^*}{d\mu}
  &= \frac{\Psi(\mu) - \kappa^{\min}(\mu)/\mu}{R\mu}
  \label{eq:dlL_dmu} \\[4pt]
  \frac{d\lambda_H^*}{d\mu}
  &= \frac{\Psi(\mu) + \kappa^{\min}(\mu)/(1-\mu)}{R(1-\mu)}
  > 0 \quad \text{always}
  \label{eq:dlH_dmu}
\end{align}
Worker $H$ always raises their upskilling fraction when equality rises
because both terms in the numerator of~(\ref{eq:dlH_dmu}) are positive.
For worker $L$, the sign depends on whether $\Psi(\mu) >
\kappa^{\min}(\mu)/\mu$, which holds when the subsistence requirement
is tight relative to current investment.

The two upskilling totals and their responses to $\mu$:
\begin{align}
  K^*_{\mathrm{ind}}(\mu) &= 2\kappa^{\min}(\mu),
  &\quad \frac{dK^*_{\mathrm{ind}}}{d\mu} &= 2\Psi(\mu) > 0
  \label{eq:dK_ind} \\[4pt]
  K^*_{\mathrm{agg}}(\mu) &= (W+1)\kappa^{\min}(\mu),
  &\quad \frac{dK^*_{\mathrm{agg}}}{d\mu} &= (W+1)\Psi(\mu) > 0
  \label{eq:dK_agg}
\end{align}
The aggregate measure responds more strongly to $\mu$ by the factor
$(W+1)/2$:
\begin{equation}
  \frac{dK^*_{\mathrm{agg}}/d\mu}{dK^*_{\mathrm{ind}}/d\mu}
  = \frac{W+1}{2}
  \label{eq:K_ratio}
\end{equation}

\begin{proposition}[Two-level upskilling neutrality]
\label{prop:neutrality}
At the constrained equilibrium:
\begin{enumerate}
  \item $K^*_{\mathrm{ind}} = 2\kappa^{\min}(\mu)$ is independent of
        $W$. Redistribution holding $\mu$ fixed does not alter the sum
        of individual investments.
  \item $K^*_{\mathrm{agg}} = (W+1)\kappa^{\min}(\mu)$ is increasing
        in $W$. A larger lower-rank population amplifies the economy-wide
        human capital base.
  \item Only $K^*_{\mathrm{agg}}$ enters firm growth through
        $\Phi(\mu) = K^*_{\mathrm{agg}}/R$.
\end{enumerate}
\end{proposition}

\subsection{Lottery Stakes, Net Transfer, and Risk Premia}

The lottery stakes respond to $\mu$ as:
\begin{equation}
  \frac{da}{d\mu} = R - \Psi(\mu) > 0,
  \qquad
  \frac{dc}{d\mu} = -R - \Psi(\mu) < 0
  \label{eq:dstakes}
\end{equation}
Worker $L$'s stake rises (income rises faster than the floor);
worker $H$'s stake falls (income falls and floor rises). The net
transfer to $L$:
\begin{equation}
  \epsilon_L(\mu) = \frac{c(\mu) - Wa(\mu)}{2}
  \label{eq:transfer_mu}
\end{equation}
is decreasing in $\mu$ when $W > 1$: higher equality reduces worker
$L$'s net lottery receipt because their larger stake means they
contribute more to the pool.

The risk premium $\tilde\rho_s = \theta_s\Sigma^2/4$ responds to $\mu$
through $\Sigma^2 = c^2 + W^2a^2$:
\begin{equation}
  \frac{d\Sigma^2}{d\mu}
  = 2c\frac{dc}{d\mu} + 2W^2 a\frac{da}{d\mu}
  = -2c(R+\Psi) + 2W^2 a(R-\Psi)
  \label{eq:dSigma}
\end{equation}
whose sign depends on the relative sizes of the stakes and $W$. For
the social premium, since $\theta_L + \theta_H$ is a fixed constant,
the sign of $d(\tilde\rho_L+\tilde\rho_H)/d\mu$ equals the sign of
$d\Sigma^2/d\mu$. Greater equality can increase or decrease the net
social welfare from the lottery depending on whether the gain in
worker $L$'s stake dominates or is dominated by the loss in worker
$H$'s stake.

\begin{proposition}[Worker responses to wage equality]
\label{prop:worker_responses}
\begin{enumerate}
  \item Worker $H$ always raises $\lambda_H^*$ when $\mu$ increases.
        Worker $L$ raises $\lambda_L^*$ when the participation
        constraint is sufficiently tight.
  \item $K^*_{\mathrm{ind}}$ is increasing in $\mu$ and independent
        of $W$.
  \item $K^*_{\mathrm{agg}}$ is increasing in both $\mu$ and $W$,
        with amplification factor $(W+1)/2$.
  \item The lottery is socially desirable iff $\theta_L + \theta_H > 0$,
        a condition independent of $\mu$.
\end{enumerate}
\end{proposition}

\subsection{Summary of Worker Responses}

Table~\ref{tab:worker_responses} collects all signed responses.

\begin{table}[h!]
\centering
\caption{Worker responses to an increase in wage equality ($\mu\uparrow$)}
\label{tab:worker_responses}
\renewcommand{\arraystretch}{1.2}
\begin{tabular}{lllp{5.5cm}}
\toprule
Object & Symbol & $\mu\uparrow$ & Mechanism \\
\midrule
Worker $L$ income & $M_L$ & $+$ & Direct \\
Worker $H$ income & $M_H$ & $-$ & Direct \\
Upskilling floor & $\kappa^{\min}$ & $+$ & Rising $\bar{b}_0 R\mu$ \\
Worker $L$ upskill fraction & $\lambda_L^*$ & $+$ (cond.) & Eq.~(\ref{eq:dlL_dmu}) \\
Worker $H$ upskill fraction & $\lambda_H^*$ & $+$ (always) & Both terms positive \\
Individual total & $K^*_{\mathrm{ind}}=2\kappa^{\min}$ & $+$ & Floor rises; $W$-free \\
Aggregate total & $K^*_{\mathrm{agg}}=(W{+}1)\kappa^{\min}$ & $+$ & Floor rises; $W$-amplified \\
Worker $L$ stake & $a$ & $+$ & Income outruns floor \\
Worker $H$ stake & $c$ & $-$ & Income falls, floor rises \\
Net transfer to $L$ & $\epsilon_L^*$ & $-$ & Larger $L$ stake hurts $L$ \\
$L$ risk premium & $\tilde\rho_L = \theta_L\Sigma^2/4$ & amb. & Via $\Sigma^2$ \\
$H$ risk premium & $\tilde\rho_H = \theta_H\Sigma^2/4$ & amb. & Via $\Sigma^2$ \\
Net social premium & $(\theta_L{+}\theta_H)\Sigma^2/4$ & amb. & Sign fixed by $\theta_L{+}\theta_H$ \\
\bottomrule
\end{tabular}
\end{table}

%----------------------------------------------------------------------
\section{Question 2: $\mu$ Constraints on Long-Term Growth}
\label{sec:firm}
%----------------------------------------------------------------------

\subsection{The Firm's Asset Accumulation Problem}

The firm accumulates assets $A_t$ according to:
\begin{equation}
  A_{t+1} = \left[(1-\pi-\delta_f) + \pi\Phi(\mu_t)\right]A_t
             - K(\mu_t) - C(A_t)
  \label{eq:firm_lom}
\end{equation}
where:
\begin{align}
  \Phi(\mu_t) &= \frac{K^*_{\mathrm{agg}}(\mu_t)}{R}
               = \frac{(W+1)\kappa^{\min}(\mu_t)}{R}
  \label{eq:Phi} \\[4pt]
  K(\mu_t) &= k|\mu_t - \mu_0|^\eta, \quad k,\eta > 0
  \label{eq:Kcost} \\[4pt]
  C(A_t) &= c_f A_t^\xi, \quad c_f, \xi > 0
  \label{eq:Ccost}
\end{align}
The upskilling index $\Phi(\mu_t)$ is the aggregate upskilling per unit
of total income. The connection between the two upskilling measures is:
\begin{equation}
  \Phi(\mu) = \frac{K^*_{\mathrm{agg}}}{R}
  = \frac{W+1}{2}\cdot\frac{K^*_{\mathrm{ind}}}{R}
  \label{eq:Phi_two_levels}
\end{equation}
When $W=1$ the two measures coincide. For $W>1$ the firm's return to
raising $\mu$ exceeds what a two-agent welfare analysis alone suggests.

\subsection{Stackelberg Timing}

The firm moves first by setting $\mu_t$, anticipating workers'
best-response functions $\lambda_s^*(\mu_t)$ from
equation~(\ref{eq:lambda_star}). By substituting these before
optimising, the firm's problem reduces to a single-variable
maximisation with no circularity:
\begin{equation}
  \mu_t^* = \arg\max_{\mu_t}\left\{
    \left[(1-\pi-\delta_f)
    + \frac{\pi(W+1)\kappa^{\min}(\mu_t)}{R}\right]A_t
    - k|\mu_t-\mu_0|^\eta - c_f A_t^\xi
  \right\}
  \label{eq:firm_opt}
\end{equation}

\subsection{Firm's First-Order Condition and Optimal $\mu^*$}

With $\eta = 2$ and $\Psi$ evaluated at reference $\Psi_0 \equiv
\Psi(\mu_0)$:
\begin{equation}
  \underbrace{\frac{\pi(W+1)\Psi_0}{R}\cdot A_t}_{%
    \text{marginal upskilling gain}}
  =
  \underbrace{2k(\mu_t^* - \mu_0)}_{%
    \text{marginal restructuring cost}}
  \label{eq:foc}
\end{equation}
Solving:
\begin{equation}
  \boxed{\mu_t^* = \mu_0 + \Gamma A_t,
  \qquad
  \Gamma \equiv \frac{\pi(W+1)\Psi_0}{2kR}}
  \label{eq:mu_opt}
\end{equation}
The coefficient $\Gamma$ is proportional to $(W+1)$: the firm's
incentive to equalise wages is driven by the aggregate upskilling
return, not individual worker responses.

\begin{proposition}[Firm's optimal wage inequality]
\label{prop:firm_mu}
\begin{enumerate}
  \item $\mu_t^* = \mu_0 + \Gamma A_t$: optimal equality is linear in
        asset wealth.
  \item $\mu_t^* > \mu_0$ for all $A_t > 0$: the firm always chooses
        greater equality than the market reference.
  \item $\partial\mu_t^*/\partial A_t = \Gamma > 0$: wealthier firms
        enforce greater equality.
  \item $\partial\mu_t^*/\partial W > 0$: larger lower-rank population
        increases the optimal equality through $(W+1)$ in $\Gamma$.
  \item $\partial\mu_t^*/\partial k < 0$: higher restructuring costs
        compress the deviation from $\mu_0$.
\end{enumerate}
\end{proposition}

\subsection{Three Constraints of $\mu$ on Long-Term Growth}

The asset growth rate is:
\begin{equation}
  g(\mu, A) = \pi\Phi(\mu) - \delta_f - c_f
              - \frac{k(\mu-\mu_0)^2}{A}
  \label{eq:growth_rate}
\end{equation}

\textbf{Constraint 1 --- Viability floor.} For a positive steady state
to exist, evaluating at $\mu = \mu_0$ (no restructuring cost) and
$A \to \infty$:
\begin{equation}
  \boxed{\pi(W+1)\kappa_0 > R(\delta_f + c_f)}
  \label{eq:viability}
\end{equation}
where $\kappa_0 = \kappa^{\min}(\mu_0)$. This is the fundamental
constraint: the aggregate upskilling return must exceed total running
costs. A larger lower-rank group ($W$ higher) makes viability easier
to satisfy. When this fails, no wage policy can generate positive
steady-state growth.

\textbf{Constraint 2 --- Restructuring upper bound.} For the firm to
benefit from $\mu_t > \mu_0$:
\begin{equation}
  \pi[\Phi(\mu_t)-\Phi(\mu_0)]A_t > k(\mu_t-\mu_0)^2
  \label{eq:upper_bound}
\end{equation}
Richer firms can sustain greater deviations from $\mu_0$. As
$A_t \to A^*$ this constraint binds at $\mu_t = \mu_t^*$ from
equation~(\ref{eq:mu_opt}).

\textbf{Constraint 3 --- Hard feasibility ceiling.} $\mu < \tfrac{1}{2}$
always: worker $H$ must earn more than worker $L$. As $\mu \to
\tfrac{1}{2}$ the upskilling gap closes and the lottery spread
$\Pi(\mu) \to 0$.

\subsection{Closed-Form Steady-State Equilibrium}

Setting $A_{t+1} = A_t = A^*$ and using $\mu^* - \mu_0 = \Gamma A^*$:
\begin{equation}
  A^*[\pi\Phi^* - \delta_f - c_f] = k\Gamma^2 A^{*2}
  \label{eq:ss_eq}
\end{equation}
Dividing by $A^* > 0$, evaluating $\Phi^*$ at $\kappa_0$:
\begin{equation}
  \boxed{A^* = \frac{\pi(W+1)\kappa_0/R - \delta_f - c_f}{k\Gamma^2}}
  \label{eq:Astar}
\end{equation}
Define three derived constants computed once from primitives:
\begin{align}
  \kappa_0 &= \frac{-1+\sqrt{1+4(\bar{b}_0 R\mu_0+\delta_w)}}{2}
  \label{eq:k0} \\[4pt]
  \Psi_0 &= \frac{\bar{b}_0 R}{1+2\kappa_0}
  \label{eq:Psi0} \\[4pt]
  \Gamma &= \frac{\pi(W+1)\Psi_0}{2kR}
  \label{eq:Gamma}
\end{align}
All steady-state objects follow sequentially:
\begin{alignat}{2}
  &\text{Firm:}\quad
  &A^* &= \frac{\pi(W+1)\kappa_0/R - \delta_f - c_f}{k\Gamma^2}
  \label{eq:Astar2}\\[4pt]
  &&\mu^* &= \mu_0 + \Gamma A^*
  \label{eq:mustar}\\[8pt]
  &\text{Workers:}\quad
  &\kappa^* &= \frac{-1+\sqrt{1+4(\bar{b}_0 R\mu^*+\delta_w)}}{2}
  \label{eq:kstar}\\[4pt]
  &&\lambda_L^* &= \frac{\kappa^*}{R\mu^*},\quad
  \lambda_H^* = \frac{\kappa^*}{R(1-\mu^*)}
  \label{eq:lambdas}\\[4pt]
  &&K^*_{\mathrm{ind}} &= 2\kappa^*,\quad
  K^*_{\mathrm{agg}} = (W+1)\kappa^*
  \label{eq:Kss}\\[4pt]
  &&a^* &= R\mu^*-\kappa^*,\quad
  c^* = R(1-\mu^*)-\kappa^*
  \label{eq:acss}\\[4pt]
  &&\Lambda^* &= Wa^*+c^*,\quad
  \epsilon_L^* = \tfrac{c^*-Wa^*}{2}
  \label{eq:Lamss}\\[8pt]
  &\text{Risk premia:}\quad
  &\tilde\rho_s^* &= \frac{\theta_s}{4}
  \left(c^{*2} + W^2a^{*2}\right)
  \label{eq:rhostar}
\end{alignat}
Every object is an explicit function of the primitives
$(R,\mu_0,W,\bar{b}_0,\delta_w,\pi,\delta_f,k,c_f,\theta_L,\theta_H)$.
Note that $(\theta_L,\theta_H)$ appear only in~(\ref{eq:rhostar})
and nowhere in equations~(\ref{eq:Astar2})--(\ref{eq:Lamss}).

\subsection{The Endogenous Kuznets Relationship}

At the steady state $\mu^* = \mu_0 + \Gamma A^*$, so wealthier
economies operate at a more equal wage split. Along the transitional
path, as $A_t$ grows from $A_0$ toward $A^*$, $\mu_t^* = \mu_0 +
\Gamma A_t$ rises monotonically. The income inequality ratio
$(1-\mu^*)/\mu^*$ falls as the economy grows: an endogenous Kuznets
relationship in which growth and equality are complementary, driven by
the firm's optimal wage policy.

%----------------------------------------------------------------------
\section{Full Comparative Statics}
%----------------------------------------------------------------------

Table~\ref{tab:cs} presents the complete steady-state comparative
statics. The $W$ row is the most important: $K^*_{\mathrm{ind}}$ is
invariant to $W$ (shown as $\mathbf{0}$) while $K^*_{\mathrm{agg}}$
rises with $W$ (shown as $+$). The risk premium parameters
$(\theta_L,\theta_H)$ appear only in the last column.

\begin{table}[h!]
\centering
\caption{Steady-state comparative statics for all three players}
\label{tab:cs}
\renewcommand{\arraystretch}{1.2}
\begin{tabular}{lcccccc}
\toprule
Shock ($\uparrow$)
  & $\partial A^*$
  & $\partial\mu^*$
  & $\partial\lambda_L^*$
  & $\partial\lambda_H^*$
  & $\partial K^*_{\mathrm{ind}}$
  & $\partial K^*_{\mathrm{agg}}$ \\
\midrule
$W\uparrow$           & $+$ & $+$ & amb. & $-$   & $\mathbf{0}$ & $+$ \\
$\pi\uparrow$         & $+$ & $+$ & $+$  & $+$   & $+$          & $+$ \\
$k\uparrow$           & $-$ & $-$ & $-$  & $-$   & $-$          & $-$ \\
$\bar{b}_0\uparrow$   & amb.& $+$ & $+$  & $-$   & $+$          & $+$ \\
$\delta_f\uparrow$    & $-$ & $-$ & $-$  & $-$   & $-$          & $-$ \\
$\delta_w\uparrow$    & $+$ & $+$ & $+$  & amb.  & $+$          & $+$ \\
$\mu_0\uparrow$       & $+$ & $+$ & amb. & $-$   & $+$          & $+$ \\
$R\uparrow$           & amb.& amb.& $-$  & $-$   & amb.         & amb.\\
\midrule
$\theta_L\uparrow$    & --- & --- & ---  & ---   & ---          & --- \\
$\theta_H\uparrow$    & --- & --- & ---  & ---   & ---          & --- \\
\bottomrule
\end{tabular}

\medskip
\small Notes: ``amb.''\ = ambiguous sign. ``---''\ = parameter does not
enter the allocation or accumulation problem. $\mathbf{0}$ denotes
exact invariance. $(\theta_L,\theta_H)$ affect only
$\tilde\rho_s^* = \theta_s(c^{*2}+W^2a^{*2})/4$: sign of $\partial\tilde\rho_s/\partial\theta_s > 0$,
sign of $\partial(\tilde\rho_L+\tilde\rho_H)/\partial(\theta_L+\theta_H) > 0$.
\end{table}

The $(\theta_L, \theta_H)$ rows confirm the separation property: risk
preference parameters have no effect on any allocation or accumulation
quantity. They affect only the welfare valuation of lottery
participation through $\tilde\rho_s^*$.

%----------------------------------------------------------------------
\section{Conclusion}
%----------------------------------------------------------------------

We have developed a three-player model in which wage inequality $\mu$
connects worker behaviour to firm growth, with all equilibrium objects
in closed form. Risk preferences are captured by two signed curvature
parameters $(\theta_L > 0, \theta_H < 0)$ that augment linear utility,
replacing the three-parameter Friedman-Savage specification while
retaining the essential asymmetry: the lower-income worker gains welfare
from the lottery and the higher-income worker loses.

On the first question --- how workers respond to $\mu$ --- the main
findings are: worker $H$ always raises their upskilling fraction when
equality rises; individual total upskilling $K^*_{\mathrm{ind}} =
2\kappa^{\min}(\mu)$ rises with $\mu$ but is invariant to $W$;
aggregate total upskilling $K^*_{\mathrm{agg}} = (W+1)\kappa^{\min}(\mu)$
rises with both $\mu$ and $W$ by amplification factor $(W+1)/2$; and
the lottery is socially desirable iff $\theta_L + \theta_H > 0$.

On the second question --- constraints of $\mu$ on growth --- the main
findings are: the viability condition $\pi(W+1)\kappa_0 > R(\delta_f+c_f)$
is the growth floor; conditional on viability the firm's optimal
$\mu^* = \mu_0 + \Gamma A^*$ is linear in wealth, producing an
endogenous Kuznets relationship; and $A^*$ is fully explicit in nine
primitives with risk parameters entering only through the welfare
evaluation of the lottery.

%----------------------------------------------------------------------
\appendix
\section{Complete Notation Table}
%----------------------------------------------------------------------

\begin{table}[h!]
\centering
\caption{All symbols, definitions, and domains}
\label{tab:notation}
\renewcommand{\arraystretch}{1.15}
\begin{tabular}{lp{8.2cm}l}
\toprule
Symbol & Definition & Domain \\
\midrule
\multicolumn{3}{l}{\textit{Primitives}} \\
$R$ & Total income to workers & $\mathbb{R}_{++}$ \\
$\mu$ & Income share of worker $L$ & $(0,\tfrac{1}{2})$ \\
$\mu_0$ & Reference (market) income split & $(0,\tfrac{1}{2})$ \\
$W$ & Population mass ratio $= i^*/(N-i^*)$ & $[1,\infty)$ \\
$\bar{b}_0$ & Subsistence coefficient; $\bar{b}=\bar{b}_0 R\mu$ & $\mathbb{R}_{++}$ \\
$\delta_w$ & Human capital depreciation & $[0,1)$ \\
$\gamma$ & Upskilling power ($=2$ for closed form) & $(0,1)$ \\
$\pi$ & Firm wage share of assets & $(0,1)$ \\
$\delta_f$ & Firm asset depreciation & $(0,1)$ \\
$k$ & Restructuring cost scale & $\mathbb{R}_{++}$ \\
$\eta$ & Restructuring cost power ($=2$ for closed form) & $(1,\infty)$ \\
$c_f$ & Maintenance cost scale & $\mathbb{R}_{++}$ \\
$\xi$ & Maintenance cost power & $(1,\infty)$ \\
$\theta_L$ & Curvature of $u_L$ at $\bar{b}$; $\theta_L > 0$ & $\mathbb{R}_{++}$ \\
$\theta_H$ & Curvature of $u_H$ at $\bar{b}$; $\theta_H < 0$ & $\mathbb{R}_{--}$ \\
\midrule
\multicolumn{3}{l}{\textit{Derived constants}} \\
$\kappa_0$ & $\kappa^{\min}(\mu_0)=(-1+\sqrt{1+4(\bar{b}_0R\mu_0+\delta_w)})/2$ & $\mathbb{R}_{++}$ \\
$\Psi_0$ & Sensitivity at $\mu_0$: $\bar{b}_0R/(1+2\kappa_0)$ & $\mathbb{R}_{++}$ \\
$\Gamma$ & Firm FOC coefficient: $\pi(W+1)\Psi_0/(2kR)$ & $\mathbb{R}_{++}$ \\
\midrule
\multicolumn{3}{l}{\textit{Worker objects}} \\
$M_L,M_H$ & Incomes $=R\mu,\;R(1-\mu)$ & $\mathbb{R}_{++}$ \\
$\lambda_s$ & Fraction of income to upskilling & $[0,1]$ \\
$\kappa_s$ & Upskilling investment $=\lambda_s M_s$ & $\mathbb{R}_+$ \\
$\ell_s$ & Lottery stake $=(1-\lambda_s)M_s$ & $\mathbb{R}_+$ \\
$b_s$ & Base income $=\kappa_s+\kappa_s^\gamma-\delta_w$ & $\mathbb{R}$ \\
$\bar{b}(\mu)$ & Subsistence floor $=\bar{b}_0R\mu$ & $\mathbb{R}_{++}$ \\
$\kappa^{\min}(\mu)$ & Min.\ upskilling; solves $\kappa+\kappa^\gamma=\bar{b}+\delta_w$ & $\mathbb{R}_{++}$ \\
$\Psi(\mu)$ & Sensitivity $=d\kappa^{\min}/d\mu$ & $\mathbb{R}_{++}$ \\
$\lambda_s^*$ & Optimal fraction $=\kappa^{\min}/M_s$ & $(0,1)$ \\
$a(\mu)$ & Worker $L$ stake $=R\mu-\kappa^{\min}$ & $\mathbb{R}_{++}$ \\
$c(\mu)$ & Worker $H$ stake $=R(1-\mu)-\kappa^{\min}$ & $\mathbb{R}_{++}$ \\
$\Pi(\mu)$ & Total pool $=Wa+c$ & $\mathbb{R}_{++}$ \\
$\Sigma^2$ & Squared spread $=c^2+W^2a^2$ & $\mathbb{R}_{++}$ \\
$\epsilon_L^*$ & Net transfer to $L$ $=(c-Wa)/2$ & $\mathbb{R}$ \\
$K^*_{\mathrm{ind}}$ & Individual upskilling $=2\kappa^{\min}$ & $\mathbb{R}_{++}$ \\
$K^*_{\mathrm{agg}}$ & Aggregate upskilling $=(W+1)\kappa^{\min}$ & $\mathbb{R}_{++}$ \\
$\tilde\rho_s$ & Risk premium $=\theta_s\Sigma^2/4$ & $\mathbb{R}$ \\
$V_s$ & Equilibrium utility $=\bar{b}+\epsilon_s+\tilde\rho_s$ & $\mathbb{R}$ \\
\midrule
\multicolumn{3}{l}{\textit{Firm objects}} \\
$A_t$ & Asset stock at period $t$ & $\mathbb{R}_+$ \\
$\Phi(\mu)$ & Upskilling index $=(W+1)\kappa^{\min}/R$ & $\mathbb{R}_{++}$ \\
$\mu_t^*$ & Optimal split $=\mu_0+\Gamma A_t$ & $(0,\tfrac{1}{2})$ \\
$A^*$ & Steady-state assets & $\mathbb{R}_+$ \\
$\mu^*$ & Steady-state split $=\mu_0+\Gamma A^*$ & $(0,\tfrac{1}{2})$ \\
\bottomrule
\end{tabular}
\end{table}

%----------------------------------------------------------------------
\appendix
\section{Derivation of Optimal $\lambda_s^*$}
\label{app:foc}
%----------------------------------------------------------------------

\subsection*{B.1\quad The Worker's Problem}

Worker $s$ maximises $V_s$ from equation~(\ref{eq:Vs_full}) subject to
$b_s(\lambda_s) \geq \bar{b}$. The objective is:
\begin{equation}
  V_s = \bar{b} + \epsilon_s(\lambda_s,\lambda_{-s})
      + \frac{\theta_s}{4}\Sigma^2(\lambda_s,\lambda_{-s})
  \tag{B.1}
\end{equation}

\subsection*{B.2\quad Monotonicity of $V_s$ in $\lambda_s$}

Differentiating with respect to $\lambda_s$ at the constrained
equilibrium $b_s = \bar{b}$:
\begin{equation}
  \frac{\partial V_s}{\partial\lambda_s}
  = \frac{\partial\epsilon_s}{\partial\lambda_s}
  + \frac{\theta_s}{4}\frac{\partial\Sigma^2}{\partial\lambda_s}
  \tag{B.2}
\end{equation}
Since $\epsilon_s = (c - W_s\ell_s - \Pi_{-s})/2$ evaluated at stakes
and $\Sigma^2 = \Pi_{-s}^2 + W_s^2\ell_s^2$, both terms have
contributions that, when combined with the base income channel
$\partial b_s/\partial\lambda_s = M_s(1+\gamma\kappa_s^{\gamma-1}) > 0$,
make $\partial V_s/\partial\lambda_s > 0$ throughout the feasible
range. This is because raising $\lambda_s$ simultaneously improves
both the win outcome (through higher $b_s$) and the lose outcome
(through higher $b_s$ and lower stake $\ell_s$). There is no interior
zero.

\subsection*{B.3\quad The Constraint Binds}

Since $V_s$ is increasing in $\lambda_s$, the worker pushes $\lambda_s$
as high as possible. The participation constraint $b_s \geq \bar{b}$
prevents $\lambda_s \to 1$ (which would leave zero lottery stake but
may still satisfy the floor). The binding condition is:
\begin{equation}
  \kappa_s^* + (\kappa_s^*)^\gamma = \bar{b} + \delta_w
  \implies
  \kappa^{\min}(\mu) = \frac{-1+\sqrt{1+4(\bar{b}+\delta_w)}}{2}
  \quad [\gamma=2]
  \tag{B.3}
\end{equation}
Since $\kappa_s^* = \lambda_s^* M_s$:
\begin{equation}
  \lambda_s^* = \frac{\kappa^{\min}(\mu)}{M_s}
  = \frac{\kappa^{\min}(\mu)}{RM_s^0}
  \tag{B.4}
\end{equation}

\subsection*{B.4\quad Ordering and Separation of Risk Preferences}

The ordering $\lambda_L^* > \lambda_H^*$ follows from $R\mu < R(1-\mu)$.
This is an income-scaling effect: the same floor $\kappa^{\min}$
represents a larger fraction of a smaller income.

The curvature parameters $(\theta_L, \theta_H)$ do not appear in
$\kappa^{\min}(\mu)$ or in $\lambda_s^*$. They enter only the risk
premium $\tilde\rho_s = \theta_s\Sigma^2/4$ in the utility
evaluation~(\ref{eq:V_final}). The allocation problem and the
preference problem are completely separated.

\end{document}
